% !TeX program = xelatex
\documentclass[12pt,a4paper]{article}

% =========================================================
% ЯЗЫК И ШРИФТЫ
% =========================================================

\usepackage{fontspec}
\usepackage{polyglossia}
\usepackage{dsfont}

\setdefaultlanguage{russian}
\setotherlanguage{english}

\setmainfont{CMU Serif}[ItalicFont=cmunti.otf]
\setsansfont{CMU Sans Serif}
\setmonofont{CMU Typewriter Text}


% =========================================================
% МАТЕМАТИКА
% =========================================================
\usepackage{amsmath}
\usepackage{amssymb}
\usepackage{amsfonts}
\usepackage{amsthm}
\usepackage{mathtools}
\usepackage{bm}
\usepackage{mathrsfs}

% =========================================================
% СТРАНИЦА
% =========================================================
\usepackage[a4paper, margin=2cm]{geometry}

% =========================================================
% ЦВЕТА И РАМКИ
% =========================================================
\usepackage{xcolor}
\usepackage[most]{tcolorbox}

% =========================================================
% ГРАФИКА
% =========================================================
\usepackage{graphicx}

% =========================================================
% СПИСКИ
% =========================================================
\usepackage{enumitem}

% =========================================================
% ЗАГОЛОВКИ
% =========================================================
\usepackage{titlesec}

% =========================================================
% КОЛОНТИТУЛЫ
% =========================================================
\usepackage{fancyhdr}

% =========================================================
% ТАБЛИЦЫ
% =========================================================
\usepackage{booktabs}
\usepackage{array}

% =========================================================
% ССЫЛКИ
% =========================================================
\usepackage{hyperref}

\hypersetup{
    colorlinks=true,
    linkcolor=blue,
    urlcolor=blue,
    citecolor=blue
}

% =========================================================
% НАСТРОЙКИ
% =========================================================
\setlength{\parindent}{0pt}
\setlength{\parskip}{6pt}

% -------------------------
% Определение
% -------------------------
\newtcolorbox{definition}[1][]{
    breakable,
    colback=red!5,
    colframe=red!75!black,
    fonttitle=\bfseries,
    title=Определение,
    #1
}

% -------------------------
% Теорема
% -------------------------
\newtcolorbox{theorem}[2][]{
    breakable,
    colback=green!5,
    colframe=green!50!black,
    fonttitle=\bfseries,
    title=Теорема #2,
    #1
}

% -------------------------
% Лемма
% -------------------------
\newtcolorbox{lemma}[2][]{
    breakable,
    colback=orange!5,
    colframe=orange!70!black,
    fonttitle=\bfseries,
    title=Лемма #2,
    #1
}

% -------------------------
% Следствие
% -------------------------
\newtcolorbox{corollary}[1][]{
    breakable,
    colback=green!5,
    colframe=green!35!black,
    fonttitle=\bfseries,
    title=Следствие,
    #1
}

% -------------------------
% Доказательство
% -------------------------
\newtcolorbox{proofbox}[2][]{
    breakable,
    colback=blue!5,
    colframe=blue!50,
    fonttitle=\bfseries,
    title=Доказательство #2,
    #1
}

% -------------------------
% Замечание
% -------------------------
\newtcolorbox{note}[1][]{
    breakable,
    colback=purple!5,
    colframe=purple!80!black,
    fonttitle=\bfseries,
    title=Замечание,
    #1
}

% -------------------------
% Свойства
% -------------------------
\newtcolorbox{props}[2][]{
    breakable,
    colback=green!5,
    colframe=green!70!black,
    fonttitle=\bfseries,
    title=Свойства #2,
    #1
}

% -------------------------
% Примеры
% -------------------------
\newtcolorbox{example}[1][]{
    breakable,
    colback=brown!5,
    colframe=brown!40!black,
    fonttitle=\bfseries,
    title=Примеры,
    #1
}

% =========================================================
% ДОКУМЕНТ
% =========================================================
\begin{document}

\tableofcontents

\newpage

\section{Линейное программирование}

\subsection{Классификация}

\textbf{По переменным:}

ограниченная --- многомерная

непрерывная --- дискретная --- смешанная

\textbf{По ограничениям:}

безусловная --- условная

\textbf{По структуре:}

глобальная --- локальная

централизованная --- распределённая --- децентрализованная

онлайн --- офлайн

статическая --- динамическая

одноуровневая --- двухуровневая --- многоуровневая, например $\min f_i(w(n), n),\ i=\overline{1,n}$,

малоразмерная --- крупномасштабная

\textbf{По целевой функции:}

стохастическая --- детерминированная

выпуклая --- невыпуклая

гладкая --- негладкая

линейная --- нелинейная

одноцелевая --- многоцелевая

\subsection{Выпуклая оптимизация}

\[
\begin{cases}
f_0(x)=x_1^2+x_2^2\to\min,\\
f_1(x)=\dfrac{x_1}{1+x_2^2}\leq 0,\\
h_1(x)=(x_1+x_2)^2=0
\end{cases}
\quad\Longrightarrow\quad
\begin{cases}
f_0(x)=x_1^2+x_2^2\to\min,\\
f_1(x)=x_1\leq 0,\\
h_1(x)=x_1+x_2=0.
\end{cases}
\]

\subsection{Задача линейного программирования}

\[
c^Tx=c_1x_1+\ldots+c_nx_n\to\max,
\]

\[
\begin{cases}
Ax\leq b,\\
x\geq 0.
\end{cases}
\]

Неравенство $\longrightarrow$ равенство.

\subsection{Каноническая форма}

\[
\begin{cases}
c^Tx\to\max,\\
Ax=b,\\
x\geq 0.
\end{cases}
\]

\[
\left[
\begin{array}{cccc|cccc}
a_{11}&a_{12}&\cdots&a_{1n}&1&0&\cdots&0\\
a_{21}&a_{22}&\cdots&a_{2n}&0&1&\cdots&0\\
\vdots&\vdots&\ddots&\vdots&\vdots&\vdots&\ddots&\vdots\\
a_{m1}&a_{m2}&\cdots&a_{mn}&0&0&\cdots&1
\end{array}
\right]
\begin{pmatrix}
x_1\\
\vdots\\
x_n\\
x_{n+1}\\
\vdots\\
x_{n+m}
\end{pmatrix}
=
\begin{pmatrix}
b_1\\
b_2\\
\vdots\\
b_m
\end{pmatrix}.
\]

\newpage

\section{Теория двойственности}

\begin{definition}
\textbf{Двойственная задача} - это задача оптимизации, которая строится по исходной (прямой) задаче так, что ограничения прямой задачи соответствуют переменным двойственной, а её целевая функция даёт оценку целевой функции прямой задачи.

Например, для прямой задачи
\[
\max c^T x,\qquad Ax\le b,\qquad x\ge0
\]
двойственная имеет вид
\[
\min b^T y,\qquad A^T y\ge c,\qquad y\ge0.
\]
\end{definition}

\begin{theorem}{1}

Задача двойственная к двойственной является прямой ЛП.

\end{theorem}

\begin{proofbox}{}
Если $x$ и $y$ --- допустимые решения прямой и двойственной задач ЛП, то
\[
c^Tx\leq b^Ty.
\]

Если
\[
c^Tx=b^Ty,
\]
то это решения оптимальные для своих задач.

\[
\begin{array}{c@{\qquad\Longleftrightarrow\qquad}c}
\begin{cases}
\min b^Ty,\\
A^Ty\geq c,\\
y\geq 0
\end{cases}
&
\begin{cases}
\max (-b)^Ty,\\
-A^Ty\leq -c,\\
y\geq 0
\end{cases}
\end{array}
\]

получаем
\[
\begin{cases}
\min -c^Tx,\\
-(A^T)^Tx\geq (-b^T)^T,\\
x\geq 0
\end{cases}
\quad\Longrightarrow\quad
\begin{cases}
\max c^Tx,\\
Ax\leq b,\\
x\geq 0.
\end{cases}
\]
\end{proofbox}

\begin{theorem}{2}

Пусть существуют допустимые $x,y$. Тогда $c^Tx\leq b^Ty$

И если $c^Tx = b^Ty$, то эти решения оптимальны для своих задач.

\end{theorem}

\begin{proofbox}{}

\[
Ax\leq b
\quad\Longrightarrow\quad
y^TAx\leq y^Tb,
\]
так как $y\geq 0$.

\[
A^Ty\geq c
\quad\Longrightarrow\quad
x^TA^Ty\geq x^Tc,
\]
так как $x\geq 0$.

Следовательно,
\[
c^Tx\leq y^TAx\leq b^Ty.
\]

\end{proofbox}

\begin{theorem}{3}

Существует $x$ --- оптимальное решение тогда и только тогда, когда существует $y$ --- оптимальное решение двойственной задачи, и
\[
\max c^Tx=\min b^Ty.
\]

\end{theorem}

\begin{lemma}{Фаркаша}

$Ax\leq b$  имеет решение $x\geq 0$ либо $A^Ty\geq 0, \ b^Ty<0$ имеет единственное решение $y\geq 0$.

\end{lemma}

\begin{proofbox}{}

\[
\begin{array}{c@{\qquad}c}
\text{прямая:} 
\begin{cases}
\max c^Tx,\\
Ax\leq b,\\
x\geq 0
\end{cases}
\text{двойственная:}
\begin{cases}
\min b^Ty,\\
A^Ty\geq c,\\
y\geq 0
\end{cases}
\end{array}
\]

\[
\begin{cases}
Ax\leq b,\\
-A^Tv\leq -c,\\
-c^Tx+b^Ty=0
\end{cases}
\Longrightarrow
\left[
\begin{array}{cc}
A&0\\
0&-A^T\\
-c^T&b^T
\end{array}
\right]
\begin{pmatrix}
x\\
y\\
\lambda
\end{pmatrix}
\leq
\begin{pmatrix}
b\\
-c\\
0
\end{pmatrix}.
\]

\[
\left[
\begin{array}{ccc}
A^T&0&-c\\
0&-A&b
\end{array}
\right]
\begin{pmatrix}
u\\
v\\
\lambda
\end{pmatrix}
\geq 0,
\]
\[
\begin{bmatrix}
b^T&-c^T&0
\end{bmatrix}
\begin{pmatrix}
u\\
v\\
\lambda
\end{pmatrix}
<0,
\qquad
u,v,\lambda\geq0.
\]

То есть
\[
A^Tu-c\lambda\geq0,
\]
\[
-Av+b\lambda\geq0,
\]
\[
b^Tu-c^Tv<0.
\]

1) Пусть $\lambda>0$:
\[
A^T\left(\frac{u}{\lambda}\right)\geq c,
\qquad
A\left(\frac{v}{\lambda}\right)\leq b,
\]
\[
b^T\left(\frac{u}{\lambda}\right)
-c^T\left(\frac{v}{\lambda}\right)<0,
\]
что противоречит теореме 2.

\medskip

2) Пусть $\lambda=0$:
\[
A^Tu\geq0,
\qquad
-Av\geq0,
\qquad
b^Tu<c^Tv.
\]

Предположим
\[
b^Tu<0. \tag{*}
\]

Так как
\[
Ax\leq b,
\]
то
\[
u^TAx\leq u^Tb.
\]
Отсюда
\[
x^TA^Tu\leq b^Tu<0,
\]
противоречие, так как
\[
x\geq0,\qquad A^Tu\geq0.
\]

Предположим
\[
c^Tv>0. \tag{**}
\]

Аналогично получаем противоречие.
\end{proofbox}

\newpage

\section{Симплекс-метод}

$P=\{x\in\mathbb{R}^n:Ax\leq b\}$ --- \textbf{полиэдр}.

\begin{definition}

\textbf{Выпуклость}

\[
\forall x^{(1)},x^{(2)}\in P
\qquad
x=\lambda x^{(1)}+(1-\lambda)x^{(2)},\quad \lambda\in[0,1] \ \Longrightarrow x\in P.
\]
\end{definition}

Пусть $P=\langle v_1,\ldots,v_k\rangle$ --- вершины

\[
x=\sum_{i=1}^k\lambda_i v_i,
\qquad
\sum_{i=1}^k\lambda_i=1.
\]

Пусть $x^*$ --- оптимальное решение задачи $\max c^Tx$

Тогда
\[
x^*=\sum_i\lambda_i v_i
\]
\[
c^Tx^*=\sum_i\lambda_i c^Tv_i
\]

Следовательно, существует $i$ такое, что $c^Tx^*\leq \max_i c^Tv_i$

Но $x^*$ --- оптимальное, поэтому $c^Tx^*=\max_i c^Tv_i$

Пусть
\[
P=\{x\in\mathbb{R}^n:Ax=b,\ x\geq0\},
\qquad
A\in\mathbb{R}^{m\times n},
\qquad
\operatorname{rank}A=m.
\]

\[
A=(B\ \overline{B}),
\qquad
x=
\begin{pmatrix}
x_B\\
x_{\overline B}
\end{pmatrix}.
\]

1) $Bx_B+\overline Bx_{\overline B}=b$.

2) $x_{\overline B}=0\Rightarrow Bx_B=b$.

$B$ --- невырождена, поэтому существует $B^{-1}$, $x_B=B^{-1}b$ --- базисное решение.

Если $x_B\geq0$, то это допустимое базисное решение.

Пусть $x_0$ --- не вершина:
\[
x=\lambda v_1+(1-\lambda)v_2,
\qquad
x_{\overline B}=0.
\]

Тогда
\[
x_B=\lambda v_{1B}+(1-\lambda)v_{2B}.
\]

\[
Bv_{1B}=b,
\qquad
Bv_{2B}=b.
\]

Следовательно,
\[
v_{1B}=\lambda B^{-1}b+(1-\lambda)B^{-1}b=B^{-1}b,
\]
откуда
\[
v_1=v_2.
\]

Возьмём $x_j\in x_{\overline B}$, $x_j>0$.

\[
Bx_B+\overline Bx_{\overline B}=b.
\]
\[
Bx_B=b-A_jx_j,
\]
\[
x_B=B^{-1}b-B^{-1}A_jx_j.
\]

Обозначим
\[
\beta=B^{-1}b,
\qquad
d=B^{-1}A_j.
\]

Тогда
\[
x_B=\beta-dx_j.
\]

\[
x_j\leq\frac{\beta}{d}.
\]

\[
\theta=\min_i\frac{\beta_i}{d_i},
\qquad d_i>0.
\]

\[
x=
\begin{pmatrix}
x_j\\
x_B
\end{pmatrix},
\qquad
c=
\begin{pmatrix}
c_j\\
c_B
\end{pmatrix}.
\]

\[
f=c_B^Tx_B+c_j^Tx_j.
\]

Так как
\[
Bx_B+\overline Bx_{\overline B}=b,
\]
то
\[
x_B=B^{-1}b-B^{-1}\overline Bx_{\overline B}.
\]

Следовательно,
\[
f
=
c_B^T
\left(
B^{-1}b-B^{-1}\overline Bx_{\overline B}
\right)
+c_{\overline B}^Tx_{\overline B},
\]
то есть
\[
f
=
c_B^TB^{-1}b+
\left(
c_{\overline B}^T-c_B^TB^{-1}\overline B
\right)x_{\overline B}.
\]

\[
f\to\max.
\]

\[
\Delta_j>0\Rightarrow x_j\uparrow,
\]
\[
\Delta_j<0\Rightarrow x_j\downarrow,
\]
\[
\Delta_j=0\Rightarrow x_j=\mathrm{const}.
\]

\newpage

\section{Решение}

\[
c^Tx\to\min,
\]
\[
Ax=b,
\]
\[
x\geq0.
\]

\[
A=[M\ N],
\]
где $M$ соответствует базисным переменным, $N$ --- небазисным.

\begin{theorem}{1}

Существует $x_0=x_0(M)$ --- допустимое базисное решение.

\[
x_0=(?, ?, ?, 0,0,0).
\]

Пусть $N_+(x_0)$ --- положительные компоненты.

\[
N_+(x_0)\subset N_B\subset N.
\]

\[
|N_B|=|M|
\quad\Longrightarrow\quad
\exists B.
\]

\end{theorem}

\begin{proofbox}{}

Пусть $B$ --- базис.

\[
A[M\,N_B]
\qquad
A[M\,N \backslash N_B]
\]

Из $A^Ty=c$ получаем $y_0[M]=c[M_B]B^{-1}$

\[
\Delta[N\backslash N_B]
=
c[M\,N_B]
-c[N_B]B^{-1}A[M\,N \backslash N_B]
\]

\[
D_B(N_B,N\backslash N_B)
=
B^{-1}A[M\,N\backslash N_B]
\]

Следовательно, $D_B(MN_B, MN_B)=E$

\[
c[N]D_B[N N \backslash N_B]
=
c[N_B]D_B(N_B,N \backslash N_B).
\]

\[
=
c[N_B]B^{-1}A[N N \backslash N_B]-c[N\backslash N_B].
\]

\[
c[N]\cdot D_B=-\Delta_0.
\]

\end{proofbox}

\begin{theorem}{2}

Допустимые решения ЗЛП имеют вид
\[
x[N]
=
x_0[N]-D_B^B x[N \backslash N_B]
\]

\end{theorem}

\begin{proofbox}{}

Пусть
\[
x=
\begin{pmatrix}
x_B\\
x_{\overline B}
\end{pmatrix}.
\]

1) Небазисные $N \backslash N_B$:
\[
x[N \backslash N_B]
=
x_0[N \backslash N_B]
-
D_B[N \backslash N_B,N \backslash N_B]x[N \backslash N_B]
=
x[N \backslash N_B]
\]

2) Базисные:
\[
x[N_B]
=
x_0[N_B]
-
D_B[N_B,N \backslash N_B]x[N \backslash N_B]
\]

Так как
\[
D_B[N_B,N \backslash N_B]
=
B^{-1}A[M,N \backslash N_B],
\]
то
\[
Ax
=
A[M,N_B]x[N_B]
+
A[M,N \backslash N_B]x[N \backslash N_B].
\]

\[
=
x_0[N_B]
-
B^{-1}
\left(
Ax-A[M,N_B]x[N_B]
\right)
\]
\[
=
x_0[N_B]-B^{-1}b+B^{-1}Bx[N_B]
=
x[N_B]
\]

3)
\[
f=cx
=
c[N]x_0[N]
+
D_0[N\backslash N_B]x[N \backslash N_B]
\]

\[
x\geq0
\Rightarrow
x_B\geq0,\quad x_{\overline B}\geq0.
\]

\medskip

Если $\Delta_0\geq0$, то $
c[M]x[M]\geq c[M]x_0[M]
\Rightarrow
x_0$ --- оптимальное

\medskip

Если $\Delta_0>0$, то
$x[N \backslash N_B]>0
\Rightarrow
cx>cx_0
\Rightarrow
\exists!\ x_0$ --- оптимальное

\end{proofbox}

\newpage

\section{Нелинейная оптимизация}

\[
\min f(x), \ x\in\mathbb{R}^n \ - \text{выпуклое} 
\]

\[
g_i(x)\leq 0, \
h_j(x)=0.
\]

\[
\Omega=\left\{x\in\mathbb{R}^n\mid g_i(x)\leq 0,\ h_j(x)=0\right\}
\]
--- множество решений.

В точке $x$ ограничения регулярны, если векторы
\[
\nabla g_i(x)\quad\text{и}\quad \nabla h_j(x)
\]
линейно независимы.

\subsection{Функция Лагранжа}

\[
L(x,\lambda,\nu)
=
f(x)
+
\sum_{i=1}^{m}\lambda_i g_i(x)
+
\sum_{j=1}^{p}\nu_j h_j(x).
\]

\[
\lambda_i\geq 0,\qquad \nu_j\in\mathbb{R}
\]
--- множители Лагранжа.

\subsection{Условия Лагранжа}

\[
\nabla_x L(x^*,\lambda^*,\nu^*)=0,
\]
\[
g_i(x^*)\leq 0,
\]
\[
h_j(x^*)=0,
\]
\[
\lambda_i^*\geq 0,
\]
\[
\lambda_i^*g_i(x^*)=0.
\]

$x^*$ --- оптимальное решение.

\subsection{Теорема Каруша--Куна--Таккера}

\textbf{Необходимые условия:}

Пусть существует $x^*$ --- локальный минимум и выполняются условия регулярности.

Тогда выполняются условия ККТ.

\textbf{Достаточные условия:}

Пусть $f(x)$ и $g_i(x)$ выпуклые, $h_j(x)$ аффинные.

Тогда
\[
\forall x^*:\ \exists \lambda^*,\nu^*
\]
удовлетворяющие условиям ККТ
\[
\Longrightarrow
\]
$x^*$ --- глобальный минимум.

\textbf{Условие Слейтера:}

Существует $x$:
\[
g_i(x)<0,\qquad h_j(x)=0.
\]

\textbf{Следствие:}

$x^*$ --- глобальный минимум
\[
\Longleftrightarrow
\exists \lambda^*,\nu^*.
\]

\subsection{Метод опорных векторов}

\[
\min \frac{1}{2}\|w\|^2,
\]
\[
y_i(w^Tx_i+b)\geq 1,\qquad i=\overline{1,N}.
\]

\[
D=\{(x_i,y_i)\}_{i=1}^N,
\qquad
y_i\in\{-1,1\}.
\]

Где x - признаки объекта, y - метка класса

\[
f(x)=w^Tx+b,
\qquad
w\in\mathbb{R}^n,\quad b\in\mathbb{R}
\]

\[
\hat y(x)=\operatorname{sign}(w^Tx+b)
\]

$w$ --- весовой вектор

\[
d_i(x,H)=\frac{|w^Tx_i+b|}{\|w\|}.
\]

Нормируем:
\[
w^Tx_1+b=1
\]
\[
w^Tx_2+b=-1
\]

\[
d=\frac{2}{\|w\|}\to\max
\]
или
\[
\frac12\|w\|^2\to\min
\]

Так как $y_i(w^Tx_i+b)=1$, то $y_i(w^Tx_i+b)\geq1$

\subsection{Решение метода опорных векторов}

\[
\min \frac{1}{2}\|w\|^2
\]

\[
1-y_i(w^Tx_i+b)\leq 0
\]

\[
\lambda_i\geq 0,\qquad \lambda_i\in\mathbb{R}
\]

\[
L(w,b,\lambda)
=
\frac{1}{2}\|w\|^2
+
\sum_{i=1}^{N}\lambda_i\left(1-y_i(w^Tx_i+b)\right)
\]

\[
=
\frac{1}{2}\|w\|^2
+
\sum_{i=1}^{N}\lambda_i
-
\sum_{i=1}^{N}\lambda_i y_i w^Tx_i
-
b\sum_{i=1}^{N}\lambda_i y_i.
\]

1.1) Стационарность по $w$

\[
\nabla_wL
=
w-\sum_{i=1}^{N}\lambda_i x_i y_i
=
0
\]

\[
w=\sum_{i=1}^{N}\lambda_i y_i x_i
\]

1.2) Стационарность по $b$

\[
\nabla_bL
=
\sum_{i=1}^{N}\lambda_i y_i
=
0.
\]

\medskip

2) $y_i(w^Tx_i+b)\geq 1$

\medskip

3) $\lambda_i\geq 0$

\medskip

4) $\lambda_i\left(1-y_i(w^Tx_i+b)\right)=0$

Пусть
\[
g=L(w,b,\lambda)\leq f(w,b)
\]

\[
L
=
\frac{1}{2}
\left(
\sum_{i=1}^{N}\lambda_i y_i x_i
\right)^2
-
\left(
\sum_{i=1}^{N}\lambda_i x_i y_i
\right)
\left(
\sum_{j=1}^{N}\lambda_j x_j y_j
\right)
+
\sum_{i=1}^{N}\lambda_i
\]

\[
=
\sum_{i=1}^{N}\lambda_i
-
\frac{1}{2}
\left(
\sum_{i=1}^{N}\lambda_i y_i x_i
\right)
\left(
\sum_{j=1}^{N}\lambda_j x_j y_j
\right).
\]

\[
\begin{cases}
\sum_{i=1}^{N}\lambda_i
-
\frac{1}{2}
\sum_{i=1}^{N}
\sum_{j=1}^{N}
\lambda_i\lambda_j x_i^Tx_jy_iy_j
\to\max\\[2mm]
\lambda_i\geq 0\\[2mm]
\displaystyle\sum_{i=1}^{N}\lambda_i y_i=0
\end{cases}
\]

\newpage

\section{Основная лемма НЛП (нелинейного программирования)}

Рассмотрим систему уравнений
\[
a_i(x)=0,
\ \
i\in I,
\ \ 
I=\{i\mid g_i(x)=0,\ i\in\{1,\ldots,p\}\}.
\]

Тогда для $g_0\in\mathbb{R}^n$, если:

1) регулярность $\nabla a_i(x)$;

2) ортогональность $g_0$:
\[
\langle\nabla a_i(x_0),g_0\rangle=0,
\qquad i\in I,
\]

можно построить параметризацию кривой $x(t)$, которая проходит в окрестности точки $t=0$, такую что
\[
x(0)=x_0
\]
\[
x'(0)=g_0
\]
\[
a_i(x(t))=0,
\ \
\forall i\in I,\quad
t\in(-\delta,\delta),\quad \delta>0
\]

Пусть
\[
a=x_1+x_2+x_3-1
\]
\[
a=x_1^2+x_2^2+x_3^2-1
\]

Пусть $x=x(t)$ --- параметризация в окрестности точки t

\[
\frac{d}{dt}a_i(x(t))
=
\left\langle
a_i'(x(t)),x'(t)
\right\rangle
\]

По теореме среднем Лагранжа:
пусть
\[
\varphi:[a,b]\to\mathbb{R},
\]
непрерывна на $[a,b]$, дифференцируема на $(a,b)$.

Тогда существует $\xi\in(a,b)$:
\[
\varphi(a)-\varphi(b)
=
\varphi'(\xi)(b-a).
\]

Для $t\neq0$:
\[
a_i(x(t))-a_i(x(0))
=
\left\langle
a_i'(x(\xi)),x'(\xi)
\right\rangle
\]

При $t\to0$
\[
\left\langle
a_i'(x(t)),x'(t)
\right\rangle=0.
\]

Так как $x(0)=x_0$, то $x'(0)=g_0, A(x)x'=0$

Условие регулярности $\Longrightarrow$ линейная независимость

Пусть
\[
J\subset N,
\qquad
|J|=|I|.
\]

Условие ортогональности:
\[
A'(x_0)g_0=0.
\]

Разобьём координаты:
\[
A'(x_0)[I,J]\,g_0[J]
+
A'(x_0)[I,N\setminus J]\,g_0[N\setminus J]
=
0.
\]

Отсюда
\[
g_0[J]
=
-
\left(
A'(x_0)[I,J]
\right)^{-1}
A'(x_0)[I,N\setminus J]
g_0[N\setminus J].
\]

Аналогично
\[
x'[J]
=
-
\left(
A'(x)[I,J]
\right)^{-1}
A'(x)[I,N\setminus J]
x'[N\setminus J].
\]

С другой стороны,
\[
x'(t)[N\setminus J]
=
g_0[N\setminus J].
\]

\[
x'[J]
=
-
\left(
A'(x)[I,J]
\right)^{-1}
A'(x)[I,N\setminus J]
g_0[N\setminus J],
\]
\[
x(0)=x_0.
\]

\begin{theorem}{}
Пусть существует система дифференциальных уравнений в нормальной форме
\[
x'(t)=F(x(t)),
\qquad
F:\mathbb{R}^n\to\mathbb{R},
\]
\[
x(0)=x_0.
\]

Тогда существует хотя бы одно решение $x(t)$ в окрестности $t=0$.

Матрица $A'(x)$ непрерывна по условию

\[
A^{-1}
=
\frac{C^T}{\det A},
\qquad
\det A\neq0,
\]
где
\[
C_{ij}=(-1)^{i+j}M_{ij}.
\]

Следовательно, $A^{-1}$ непрерывна
\end{theorem}

\begin{corollary}
Cуществует $x(t)$.

При $t=0$:
\[
x'(0)[N\setminus J]=g_0[N\setminus J],
\]
\[
x'(0)[J]
=
-
\left(
A'(x_0)[I,J]
\right)^{-1}
A'(x_0)[I,N\setminus J]
g_0[N\setminus J].
\]

Но
\[
g_0[J]
=
-
\left(
A'(x_0)[I,J]
\right)^{-1}
A'(x_0)[I,N\setminus J]
g_0[N\setminus J].
\]

Следовательно,
\[
g_0[J]=x'(0)[J],
\]
то есть
\[
x'(0)=g_0.
\]

\end{corollary}

\end{document}